% =============================================================================
% Benötigte Pakete: inputenc, fontenc, babel, amsmath, booktabs,
%                   graphicx, hyperref, biblatex, siunitx, acronym
% =============================================================================

\chapter{Grundlagen}
\label{chap:grundlagen}
% TODO: update this when done proofreading
Dieses Kapitel legt das theoretische Fundament für die in den folgenden Kapiteln
beschriebene Entwicklung. Es gliedert sich in fünf Abschnitte: Zunächst werden die
physikalischen Grundlagen der LoRa"=Modulationstechnologie erläutert, gefolgt von
einer Darstellung des LoRaWAN"=Protokollstacks und seiner Netzwerkarchitektur.
Anschließend werden Sicherheitsmechanismen, Energieeffizienz sowie Methoden der
RF"=Planung und Signalausbreitung behandelt. Abschließend wird ein kurzer Überblick
über die Normierungsanforderungen im Bereich Schutzsensorik gegeben.

% -----------------------------------------------------------------------------
\section{LoRa-Modulation}
\label{sec:lora-physical}
% -----------------------------------------------------------------------------
\subsection{Allgemeines}

\acrshort{lora} ist eine von der Firma Semtech entwickelte Modulationstechnologie und bildet die physikalische Grundlage des 
\acrshort{lorawan}"=Protokolls.  Es basiert auf der \acrfull{css}"=Modulation.
bei der das zu übertragende Signal auf ein breiteres Frequenzband gespreizt wird, indem
sogenannte Chirps-Signale mit kontinuierlich steigender oder fallender Frequenz
als Trägersymbole verwendet. % TODO: \cite{augustin2016study}. add pics of chirps

Ursprünglich für Militäranwendungen gedacht, hat das CSS-Verfahren in den letzten 20 Jahren in der Nachrichtentechnik 
zunehmend an Bedeutung gewonnen dank seiner hervorragenden Eigenschaften \cite{LoRaModulationBasics}: 

\paragraph*{Niedriger Energieverbrauch}
Ähnlich wie \acrshort{fsk} besitzt LoRa eine konstante Einhüllende, wodurch kostengünstige, energieeffiziente
Leistungsverstärker verwendet werden können.
\paragraph*{Hohe Störfestigkeit} 
Da ein LoRa-Symbol lang und "ausgespreizt" ist, haben kurz Rauschimpulse nur einen verschwindenen Einfluss auf die Demodulation. 
Aus diesem Grund sind LoRa-Signale sehr resistent gegenüber In-Band- und Out-of-Band-Interferenzen.
Eine Gleichkanalsunterdrückung von über \SI{20}{\dB} kann erreicht werden, verglichen mit nur \SI{-6}{\dB} bei FSK.
\paragraph*{Mehrwegeausbreitung und Fading}
Das breitbandige Chirp-Signal bietet eine hohe Immunität gegenüber Mehrwegeausbreitung und Fading, was den
Einsatz in urbanen und suburbanen Umgebungen begünstigt.
\paragraph*{Hohe Reichweite}
Für eine gegebene Sendeleistung kann bei LoRa die Reichweite bis zu 4 mal so groß im Vergleich zu FSK.

% TODO: figure what to do with this
Gegenüber klassischen Schmalband"=Modulationsverfahren wie \acrshort{fsk} besitzt CSS einen
wesentlichen Vorteil: Das Signal kann noch korrekt demoduliert werden, selbst wenn der
Empfangspegel deutlich unterhalb des Rauschpegels liegt. Das resultierende
Signal"=zu"=Rausch"=Verhältnis (\acrshort{snr}) kann dabei negative Werte von bis zu
\qty{-20}{dB} annehmen. Mit FSK lassen sich Datenraten von bis zu \qty{50}{kbps}
erzielen; LoRa liefert demgegenüber typische Nutzdatenraten zwischen \qty{300}{bps}
und \qty{11}{kbps}.

\subsection{Wichtige Kenngrößen und Formeln} % TODO: work on this

\subsection{Spreading Factor, Bandbreite und Coding Rate}
\label{subsec:sf-bw-cr}

Die wichtigsten Parameter zur Steuerung des LoRa"=Signals sind der \acrfull{sf}, 
die Bandbreite (\acrshort{bw}) sowie die Coding Rate (\acrshort{cr}). % TODO: cite sylvain montagny

Der SF gibt die Anzahl der Bits pro Symbol bzw Chirp an und kann Werte von 7 bis 12 annehmen. 
Wenn zum Beispiel mit SF12 übertragen wird, enthält jeder Chirp 12 Bits. Es gibt in diesem Fall also $2^{12} = 4096$ mögliche Symbole, 
die jeweils eine eindeutige Bitfolge repräsentieren. Die Übertragungszeit pro Symbol lässt sich berechnen durch:

\begin{equation}
  T_\text{sym} = \frac{2^\text{SF}}{\text{BW}}
  \label{eq:tsym}
\end{equation}

% TODO: check this
Die Bandbreite (\acrshort{bw}) kann 125, 250 oder \qty{500}{kHz} betragen. Eine größere Bandbreite ermöglicht höhere 
Datenraten, reduziert jedoch die Empfindlichkeit des Empfängers und damit die Reichweite.

Die \acrlong{toa} bzw. Gesamtübertragungszeit eines LoRa-Pakets folgt unmittelbar:

\begin{equation}
  T = n_{sym} \cdot T_\text{sym} = n_{sym} \cdot \frac{2^\text{SF}}{\text{BW}}
  \label{eq:ToA}
\end{equation}

Der SF bestimmt, wie stark das Signal gespreizt wird, und kann ganzzahlige Werte
zwischen 7 und 12 annehmen. Ein höherer SF erhöht die Empfindlichkeit des
Empfängers und damit die erreichbare Reichweite, verlängert jedoch gleichzeitig die
Übertragungszeit eines Paketes erheblich und erhöht den
Energieverbrauch pro Übertragung. Tabelle~\ref{tab:sf_vergleich} fasst die
Auswirkungen der verschiedenen \acrshort{sf}s zusammen \cite{LoRaModulationBasics}.

\begin{table}[htbp]
  \centering
  \caption{Vergleich der LoRa Spreading Factors bei \qty{125}{kHz} Bandbreite}
  \label{tab:sf_vergleich}
  \begin{tabular}{@{}lrrrl@{}}
    \toprule
    \textbf{SF} & \textbf{Datenrate} & \textbf{Empfindlichkeit} &
    \textbf{Time on Air} & \textbf{Typischer Einsatz} \\
                & [bps]        & [dBm]              &
    (20\,Byte) [ms] & \\
    \midrule
    SF7  & 5470 & $-123$ &  41 & Kurze Distanz, Energie sparen \\
    SF8  & 3125 & $-126$ &  72 & Ausgewogener Betrieb \\
    SF9  & 1760 & $-129$ & 144 & Mittlere Distanz \\
    SF10 &  980 & $-132$ & 289 & Grenzbereich Stadtgebiet \\
    SF11 &  440 & $-134$ & 578 & Ländliche Umgebung \\
    SF12 &  250 & $-137$ & 991 & Maximale Reichweite \\
    \bottomrule
  \end{tabular}
\end{table}

Als Faustregel gilt: SF7 zusammen mit einer Bandbreite von \qty{125}{kHz}
(\texttt{SF7BW125}) stellt den optimalen Ausgangspunkt für die meisten Anwendungen
dar, da die Luftzeit und damit der Energieverbrauch minimiert werden. 
Die Coding Rate beschreibt den Anteil der redundanten Bits zur Fehlerkorrektur und
nimmt Werte von 4/5 bis 4/8 an.

\subsection{Frequenzbänder und regulatorische Rahmenbedingungen}
\label{subsec:frequenzbaender}

LoRa nutzt sogenannte \acrfull{ism}"=Bänder, für die keine Lizenz erworben werden muss.
In Europa werden die Frequenzbänder \SI{433}{\MHz} und \SI{868}{\MHz} verwendet,
in Nordamerika \SI{915}{\MHz}. Da die Frequenzzuweisungen regional verschieden sind,
können Geräte, die für den europäischen Markt entwickelt wurden, nicht ohne
Anpassungen in anderen Regionen betrieben werden~% TODO: \cite{lorawan_spec}.

Für das \SI{868}{\MHz}"=Band in Europa gilt eine Duty"=Cycle"=Begrenzung: Endgeräte
dürfen je nach Subband 0{,}1\,\%, 1\,\% oder 10\,\% der Zeit senden. Die
10\,\%"=Ausnahme gilt ausschließlich für Gateways~% TODO: \cite{lorawan_spec}. In der Praxis
bedeutet eine 1\,\%"=Beschränkung beispielsweise, dass ein Gerät nach einer
Sendedauer von \SI{1}{\s} mindestens \SI{99}{\s} warten muss, bevor es erneut auf
demselben Kanal sendet. Für zeitkritische Warn"= und Alarmsysteme ist diese
Einschränkung bei der Systemauslegung zwingend zu berücksichtigen.

Die maximale Paketgröße in Europa beträgt 222\,Bytes bei hoher Datenrate und
51\,Bytes bei niedriger Datenrate. Dabei addiert LoRaWAN stets einen
13"=Byte"=Overhead zum Nutzdaten"=Payload~% TODO: \cite{lorawan_spec}.

% -----------------------------------------------------------------------------
\section{LoRaWAN -- Network Layer} % TODO: remove --
\label{sec:lorawan-network}
% -----------------------------------------------------------------------------

Der Begriff \acrfull{lorawan} definiert den Netzwerkstandard, dem die LoRa-Modulationstechnologie zugrunde liegt.
Der LoRaWAN-Standard wird von der sogenannten LoRa-Alliance - einem im Jahre 2015 gegründeten
Bündnis internationaler Unternehmen - gepflegt und in ihren umfangreichen technischen Spezifikationen dokumentiert.
Diese beginnen alle mit dem Präfix "TS", gefolgt von einer dreistelligen Dokumentennummer. 

\subsection*{Netzwerkarchitektur}
\label{subsec:netzwerkarchitektur}

Eine LoRaWAN"=Kommunikationskette setzt sich aus drei Gliedern zusammen: Endgeräten/\acrfull{ed}, Gateways und Servern. 
Die Server"=Seite gliedert sich weiter in \acrfull{ns}, \acrfull{js} und \acrfull{as}. % TODO: add pic

\subsection{Endgeräte}

Auf der untersten Ebene befinden sich die \textbf{Endgeräte}, deren Zentraleinheit Mikrokontroller (\acrshort{mcu}) mit entsprechender
Stromversorgung darstellen. Auf dem MCU ist ein LoRaWAN"=Stack integriert, der als Schnittstelle zur LoRa-Funkeinheit fungiert.
Diese übernimmt wie bereits beschrieben die physikalische Modulation und Demodulation der Signale, welche mithilfe einer Antenne
in die Luft gesendet bzw. von der Luft empfangen werden können. Zusätzlich können Endgeräte mit anwendungsspezifischen Peripheriegeräten 
wie Sensoren oder Aktoren ausgestattet sein. 

\subsubsection*{Geräteklassen}

Das Dokument \citetitle{LoRaWANL21042020} definiert drei Geräteklassen (A, B und C), die festlegen,
wann ein Endgerät Downlink"=Nachrichten, also vom Server zu Endgerät, empfangen kann. Die Klasse kann während des
Betriebs gewechselt werden.

\textbf{Klasse A} ist die energieeffizienteste und am weitesten verbreitete Klasse, denn jedes LoRaWAN-Gerät muss sie implementieren.
Nach jeder Uplink"=Übertragung öffnet das Gerät zwei kurze Fenster (RX1 und RX2) für das Empfangen von Downlink"=Nachrichten. 
Außerhalb dieser Fenster befindet sich das Gerät im Schlafmodus, um den Stromverbrauch zu minimieren. Eine erneute Uplink-Kommunikation 
darf erst nach Empfang einer Downlink-Nachricht oder nach Ablauf von RX2 erfolgen. Diese Klasse eignet sich besonders für batteriebetriebene Sensoren, die 
nur gelegentlich Daten senden müssen, wie es bei vielen Schutzsensoren der Fall ist. % TODO: add pic on page 12 TS001

Im Vergleich zur Klasse A ist die \textbf{Klasse B} flexibler, da sie einen periodischen Empfang durch zeitgesteuerte \acrshort{rx}-Fenster ermöglicht,
die mittels Gateway"=\textbf{B}eacons synchronisiert werden. Dadurch würde unmittelbar der Network Server erfahren, sobald Endgeräte empfangsbereits sind.
Dies erhöht die Reaktionsfähigkeit auf Kosten eines höheren Energieverbrauchs. Wenn ein Gerät sich als Klasse-B-Gerät Network-Server gegenüber 
identifizieren will, muss der entsprechende \texttt{ClassB}-Bit in einem Uplink-Signal gesetzt werden. 

\textbf{Klasse C} ermöglicht einen kontinuierlichen Empfang, indem das Endgerät nur während des Sendens nicht empfängt. 
Das hat konsequenterweise den Nachteil, dass Klasse-C-Geräte für batteriebetriebene Geräte ungeeignet sind
und für das Thema Off-Grid-Schutzsensorik eine eher untergeordnete Rolle spielen.

\begin{table}[ht]
    \centering
    \caption{Zusammenfassung der LoRaWAN-Klassen und ihrer Merkmale}
    \label{tab:lorawan_mnemonics}
    \begin{tabularx}{\textwidth}{l l X X}
        \toprule
        \textbf{Klasse} & \textbf{Merkhilfe} & \textbf{Hauptmerkmal} & \textbf{Energieprofil} \\
        \midrule
        \textbf{A} & \textbf{A}ll & Grundfunktion: Muss von \textbf{allen} Geräten unterstützt werden. & \textbf{Maximal} (Schlafmodus standardmäßig) \\
        \addlinespace
        \textbf{B} & \textbf{B}eacon & Synchronisation durch Gateway-\textbf{B}eacons (erfordert \texttt{ClassB}-Bit). & \textbf{Medium} (periodisches Aufwachen) \\
        \addlinespace
        \textbf{C} & \textbf{C}ontinuous & \textbf{C}ontinuierlicher Empfang (fast immer aktiv). & \textbf{Minimal} (nicht für Batteriebetrieb) \\
        \bottomrule
    \end{tabularx}
\end{table}

\subsection{Gateways}

Nach den Endgeräten schließen befinden sich direkt im Anschluss die \textbf{Gateways}, welche ebenfalls auf der Physical layer % TODO: translate
wirken. Bei der uplinkmäßigen Übertragung, also von Endgerät zu Server, empfangen sie die Signale der Endgeräte und leiten 
diese an die Netzwerkserver weiter. Umgekehrt gilt für die Downlink-Richtung, 
dass Gateways die Signale der Server Endgeräten zur Verfügung stellen. Da Gateways permanent auf allen Kanälen lauschen, 
können Endgeräte jederzeit senden, solange regulatorische Gegebenheiten es erlauben. 

Die Kommunikation zwischen Endgeräten und Gateways basiert auf dem ALOHA"=Protokoll. Das bedeutet, dass ein Endgerät nicht an 
ein bestimmtes Gateway gebunden ist, sondern eine beliebige Anzahl an Gateways in Reichweite nutzen kann, um mit dem Netzwerk 
Daten auszutauschen.

Aufgrund dessen, dass Gateways aus Sicht des LoRaWAN"=Protokolls transparent sind, werden sie in technischen Spezifikationen der 
LoRa-Alliance nur flüchtig erwähnt und größtenteils "wegabstrahiert". Um diese theoretische Transparenz jedoch zu gewährleisten, 
müsste man bei der Platzierung Gedanken machen. Indoor-Gateways zum Beispiel, also in Innenräumen installierte Gateways, 
haben eine beschränkte Reichweite, die sie dem hohen Dämpfungsvermögen von Baumaterialien zu verdanken haben. 
Daher empfiehlt es sich, sie in der Nähe von Fensterfronten zu platzieren, um den Signaldurchgang zu verbessern. 

Für den Außeneinsatz sind hingegen Outdoor"=Gateways erforderlich, die aufgrund wasserdichter Gehäuse und Hohlraumfilter 
teurer sind als ihre Indoor"=Pendants. In diesem Fall gilt es zu beachten, dass Gateways möglichst hoch montiert werden sollen, 
um die Sichtlinie zu den Endgeräten zu garantieren und damit die Reichweite zu erhöhen.

\subsection{Network-Server}

Mit den Gateways hört auch der Geltungsbereich von LoRa auf und wir verlassen somit die Bitübertragungsschicht des OSI-Modells. % TODO: check validity of this statement
Die Informationspakete gelangen als Nächstes über eine Backhaul-Verbindung (bswp. Ethernet oder Mobilfunk) zu den Servern, 
von welchen der \textbf{Network-Server} (\acrshort{ns}) der zentrale Knotenpunkt der LoRaWAN"=Sterntopologie ist. 
An dieser Stelle tritt das Dokument \citetitle{LoRaWanBackendInterfaces2020} zum Vorschein. 

Zu den Aufgaben eines Network-Servers gehören: Überprüfung von Endgerätadressen, Paketsauthentifizierung, Anpassung der Datenraten, Weiterleiten von Payloads bzw. Join-Anfragen an entsprechende 
Application-Servers und Join-Server jeweils. Zusätzlich kümmern sich Network-Server um die Konsolidierung von Nachrichten, 
die von mehreren Gateways empfangen wurden, sowie um die Überwachung der Netzwerk- und Gateway"=Leistung. % TODO: NRM model

\subsubsection*{Adaptive Data Rate}
\label{subsec:adr}

\acrfull{adr} ist ein Mechanismus des Network Servers zur automatischen Anpassung von
Datenrate und Sendeleistung eines Endgeräts. Endgeräte in der Nähe eines Gateways
werden auf höhere Datenraten (niedrigere SF"=Werte) umgeschaltet, was die Luftzeit
verkürzt und die Netzwerkkapazität erhöht. Geräte in größerer Entfernung werden mit
niedrigeren Datenraten betrieben, um ausreichende Verbindungsqualität sicherzustellen.
Wird ein weiteres Gateway zum Netzwerk hinzugefügt, aktualisieren die Endgeräte ihren
SF automatisch. ADR ist ausschließlich für stationäre
Geräte geeignet, für mobile Anwendungen wie das hier betrachtete Schutzsensorsystem 
muss ADR deaktiviert oder durch ein geeignetes mobilitätsbewusstes Verfahren ersetzt
werden.

\subsection{Join-Server}

Der \textbf{Join-Server} verwaltet den Netzwerkbeitritt bzw. Aktivierung von Endgeräten. Jeder \acrshort{js} verfügt zur Identifikation über
eine eindeutige JoinEUI (\acrlong{eui}), welche von den Endgeräten bei Join-Anfragen verwendet wird. Erfolgreiche Beitrittsanfragen bewirken
die Erzeugung von Sitzungsschlüsseln (\acrshort{NwkSKey} und \acrshort{AppSKey}), somit ist die Anbindung von 
Endgeräten an das LoRaWAN-Netzwerk abgeschlossen und sie können mit der Datenübertragung beginnen. % TODO: maybe rewrite this

\subsubsection*{Aktivierungsverfahren: OTAA und ABP}
\label{subsec:aktivierung}

Es wird zwischen zwei Aktivierungsverfahren unterschieden: \acrfull{otaa} und \acrfull{abp}.

Bei OTAA initiiert das Gerät beim ersten Start eine \textit{Join Request}, die sowohl seine \acrshort{DevEUI}
als auch die Join-EUI des Join-Servers enthält. Zusätzlich wird noch eine DevNonce übermittelt, die bei jeder Anfrage
inkrementiert werden soll, andernfalls würde der die Aktivierung verweigert. Bei erfolgreicher Aktivierung sendet der 
Join-Server eine \textit{Join Accept} und stellt dabei die notwendigen Sitzungsschlüssel, \acrfull{NwkSKey} und \acrfull{AppSKey}, bereit.
Das Gerät bekommt außerdem eine eindeutige \acrfull{DevAddr} vom Network-Server zugewiesen. Die LoRa-Alliance empfiehlt
OTAA als bevorzugtes Aktivierungsverfahren, da die dynamische Ableitung von Sitzungsschlüsseln für Sicherheit sorgt. 
% TODO: add pic illustrating OTAA

Im Gegensatz zu OTAA werden bei \acrshort{abp} alle benötigten Parameter (\acrshort{DevEUI}, \acrshort{DevAddr} und Sitzungsschlüssel)
vorab im Gerät fest einprogrammiert. Das Gerät kann sofort ohne Join"=Prozedur senden. Dieses Verfahren ist einfacher zu implementieren, 
allerdings dafür weniger flexibel und gilt sicherheitstechnisch als suboptimal. Wenn beispielsweise ein Gerät kompromittiert wird, 
können die fest eingebauten Schlüssel nicht mehr geändert werden, was zu einem dauerhaften Sicherheitsrisiko führt.

% TODO: see where this section fits
\subsection{Öffentliche vs. private Netzwerke}
\label{subsec:netzwerktypen}

Bei der Wahl der Netzwerkinfrastruktur sind drei Faktoren maßgeblich: Kosten,
Zugangskontrolle und Datenschutz~\cite{semtech_lora_basics}.

Öffentliche Netze (z.B. The Things Network) sind für kleine Gerätezahlen und
Prototyping geeignet, erfordern jedoch eine bestehende Abdeckung und bieten keine
Möglichkeit zur Optimierung durch eigene Gateways. Private Netzwerke eignen sich
für große Gerätezahlen und sicherheitskritische Anwendungen, erfordern jedoch die
Installation und den Betrieb eigener Gateways und Server"=Infrastruktur. Für das
in dieser Arbeit betrachtete Off"=Grid"=Szenario -- den Einsatz in abgelegenen
Gefahrgebieten ohne bestehende Infrastruktur -- ist ein privates, autarkes Netzwerk
die einzig praktikable Option. Als Network"=Server"=Software kommt dabei \textbf{ChirpStack}
zum Einsatz, da dieser vollständig lokal auf einem Raspberry Pi betrieben werden kann
und keine Internetverbindung erfordert. Semtech bietet für Prototyping"=Zwecke zudem
einen kostenlosen Network Server für bis zu drei Gateways und zehn Endgeräte an.

\subsection{Einfluss des Spreading Factors auf den Energieverbrauch}
\label{subsec:sf-energie}

Die Wahl des Spreading Factors hat direkten Einfluss auf den Energieverbrauch pro
übertragenem Byte. SF12 erzielt zwar die größte Reichweite, verlängert jedoch die
Luftzeit und damit die Zeit, in der der Sender aktiv ist, um etwa das 24"=fache
gegenüber SF7 (vgl. Tabelle~\ref{tab:sf_vergleich}). Die energieoptimale Strategie
besteht darin, den niedrigstmöglichen SF zu verwenden, der noch eine zuverlässige
Übertragung gewährleistet.

\subsection{Link Budget}
\label{subsec:link-budget}

Das Link Budget beschreibt die Gesamtbilanz aller Signalgewinne und -verluste
entlang der Übertragungsstrecke. Das \textit{Maximum Allowable Path Loss} (MAPL)
ergibt sich aus:

\begin{equation}
  \text{MAPL} = P_\text{TX} + G_\text{TX} + G_\text{RX} - L_\text{cable}
                - \text{NF} - \text{SNR}_\text{min} - \text{Margin}
  \label{eq:mapl}
\end{equation}

wobei $P_\text{TX}$ die Sendeleistung in \si{dBm}, $G_\text{TX}$ und $G_\text{RX}$
die Antennengewinne in \si{dBi}, NF das Rauschen des Empfängers und
$\text{SNR}_\text{min}$ das minimale empfängerseitig tolerierte Signal"=zu"=Rausch"=Verhältnis
beschreiben. Als maximaler Antennengewinn für eine einfache Dipolantenne gilt
$G = \SI{2.15}{dBi}$. Ein Sicherheitsabstand von 5 bis \SI{10}{\dB} wird empfohlen,
um Schwundeffekte und Hindernisse zu kompensieren~\cite{semtech_lora_basics}.
Uplink und Downlink"=Budget sind aufgrund unterschiedlicher Antennenhöhen in der Regel
nicht symmetrisch.

% Format : \newacronym{label}{abbreviation}{full term}
% TODO: remember to run makeglossaries main 
\newacronym{lora}{LoRa}{Long Range}
\newacronym{lorawan}{LoRaWAN}{Long Range Wide Area Network}
\newacronym{css}{CSS}{Chirp Spread Spectrum}
\newacronym{sf}{SF}{Spreading Factor}
\newacronym{toa}{ToA}{Time-on-Air}
\newacronym{bw}{BW}{Bandwidth}
\newacronym{cr}{CR}{Coding Rate}
\newacronym{snr}{SNR}{Signal-to-Noise Ratio}
\newacronym{fsk}{FSK}{Frequency Shift Keying}
\newacronym{mcu}{MCU}{Microcontroller Unit}
\newacronym{ism}{ISM}{Industrial, Scientific, and Medical}
\newacronym{phy}{PHY}{Physical Layer (Bitübertragungsschicht)}
\newacronym{ed}{ED}{End Device}
\newacronym{ns}{NS}{Network Server}
\newacronym{js}{JS}{Join Server}
\newacronym{as}{AS}{Application Server}
\newacronym{ts}{TS}{Technical Specification}
\newacronym{rx}{RX}{Receive (Empfangen)}
\newacronym{tx}{TX}{Transmit (Senden)}
\newacronym{adr}{ADR}{Adaptive Data Rate}
\newacronym{otaa}{OTAA}{Over-the-Air Activation}
\newacronym{abp}{ABP}{Activation by Personalization}
\newacronym{DevAddr}{DevAddr}{Device Address}
\newacronym{NwkSKey}{NwkSKey}{Network Session Key}
\newacronym{AppSKey}{AppSKey}{Application Session Key}
\newacronym{NwkAddr}{NwkAddr}{Network Address}
\newacronym{eui}{EUI}{Extended Unique Identifier}
\newacronym{DevEUI}{DevEUI}{Device EUI}
\newacronym{mapl}{MAPL}{Maximum Allowable Path Loss}
\newacronym{fspl}{FSPL}{Free Space Path Loss}
\newacronym{srtm}{SRTM}{Shuttle Radar Topography Mission}
\newacronym{dguv}{DGUV}{Deutsche Gesetzliche Unfallversicherung}